\begin{figure}[H]
\centering
\includegraphics[width=0.86\textwidth]{c6_62_p450.jpg}
\caption*{\textbf{表 6.5}\ 双值点群的对称适配函数（原书第 437 页起扫描）}
\end{figure}
\begin{figure}[H]
\centering
\includegraphics[width=0.86\textwidth]{c6_62_p451.jpg}
\caption*{原书第 438 页}
\end{figure}
\begin{figure}[H]
\centering
\includegraphics[width=0.86\textwidth]{c6_62_p452.jpg}
\caption*{原书第 439 页}
\end{figure}
\begin{figure}[H]
\centering
\includegraphics[width=0.86\textwidth]{c6_62_p453.jpg}
\caption*{原书第 440 页}
\end{figure}
\begin{figure}[H]
\centering
\includegraphics[width=0.86\textwidth]{c6_62_p454.jpg}
\caption*{原书第 441 页}
\end{figure}
\begin{figure}[H]
\centering
\includegraphics[width=0.86\textwidth]{c6_62_p455.jpg}
\caption*{原书第 442 页}
\end{figure}
\begin{figure}[H]
\centering
\includegraphics[width=0.86\textwidth]{c6_62_p456.jpg}
\caption*{原书第 443 页}
\end{figure}
\begin{figure}[H]
\centering
\includegraphics[width=0.86\textwidth]{c6_62_p457.jpg}
\caption*{原书第 444 页}
\end{figure}
\begin{figure}[H]
\centering
\includegraphics[width=0.86\textwidth]{c6_62_p458.jpg}
\caption*{原书第 445 页}
\end{figure}
\begin{figure}[H]
\centering
\includegraphics[width=0.86\textwidth]{c6_62_p459.jpg}
\caption*{原书第 446 页}
\end{figure}
\begin{figure}[H]
\centering
\includegraphics[width=0.86\textwidth]{c6_62_p460.jpg}
\caption*{原书第 447 页}
\end{figure}
\begin{figure}[H]
\centering
\includegraphics[width=0.86\textwidth]{c6_62_p461.jpg}
\caption*{原书第 448 页}
\end{figure}
\begin{figure}[H]
\centering
\includegraphics[width=0.86\textwidth]{c6_62_p462.jpg}
\caption*{原书第 449 页}
\end{figure}
\begin{figure}[H]
\centering
\includegraphics[width=0.86\textwidth]{c6_62_p463.jpg}
\caption*{原书第 450 页}
\end{figure}
\begin{figure}[H]
\centering
\includegraphics[width=0.86\textwidth]{c6_62_p464.jpg}
\caption*{原书第 451 页}
\end{figure}
\begin{figure}[H]
\centering
\includegraphics[width=0.86\textwidth]{c6_62_p465.jpg}
\caption*{原书第 452 页}
\end{figure}
\begin{figure}[H]
\centering
\includegraphics[width=0.86\textwidth]{c6_62_p466.jpg}
\caption*{原书第 453 页}
\end{figure}